\subsection{Methods To Represent Sets}

There are 3 primary ways of representing sets.

\begin{enumerate}
  \item \textbf{Sentence Form} --- The elements of the set are basically described using sentences.
	\begin{itemize}
	  \item This method should be avoided if possible.
	  \item Tends to be very wordy. May lack necessary precision.
	  \item Example: Let $S$ be the set of whole numbers greater than or equal to twenty and less than or equal to 99.
	\end{itemize}

  \item \textbf{Roster Form} --- Similar to a classroom roster, it is an explicit list of the elements of a set.
	\begin{itemize}
	  \item Very easy to understand.
	  \item Typically useful for very small sets or sets with elements that form obvious patterns.
	  \item Ellipsis can be used only if an obvious pattern can be established.
		\begin{itemize}
		  \item Example: $A = \{2,3,-4,9\}$
		  \item Example: $B = \{1,2,3,\ldots,99,100\}$
		  \item Example: $\mathbb{N} = \{1,2,3,\ldots\}$
		  \item Example: $\mathbb{Z} = \{\ldots,-3,-2,-1,0,1,2,3,\ldots\}$
		  \item Bad Example: $K = \underbrace{\{1,2,5,10,17,\ldots\}}_{\text{Pattern Is Not Obvious}}$
		\end{itemize}
	  \item Deciding on whether to use roster form may depend on the mathematical maturity of your audience.
		\begin{itemize}
		  \item Example: $P = \{2,3,5,7,11,13,17,19,23,\ldots\}$ (Prime Numbers)
		  \item Example: $T = \{1,4,9,16,25,36,49,\ldots\}$ (Perfect Squares)
		\end{itemize}
	\end{itemize}

  \item \textbf{Set Builder Form} --- Requires the most mathematical maturity to use. However, it becomes second nature once familiar.
	\begin{itemize}
	  \item Resembles a food recipe. You have a list of ingredients and instructions on how to use the ingredients.
	  \item Often used when sets need to be constructed or can be described using some form of construction.
		\begin{itemize}
		  \item Example:
			\[
			  E = \{2k \;\mid\; k \in \mathbb{W}\} = \{0,2,4,6,8,10,\ldots\}
			\]
			(Set of even numbers --- whole numbers that are a multiple of two)
		  \item Example:
			\[
			  O = \{2k+1 \;\mid\; k \in \mathbb{W}\} = \{1,3,5,7,9,\ldots\}
			\]
			(Set of odd numbers --- whole numbers that are one plus a multiple of two)
		\end{itemize}
	\end{itemize}
\end{enumerate}

